\chapter{LITERATURE SURVEY}

\section{Overview of Bus Ticket Booking System}

The Online Bus Booking System is a web-based application designed to simplify the process of booking bus tickets. It allows users to search buses by route and date, view available schedules, select seats, and complete bookings efficiently.

The system also includes user authentication and an admin module to manage bus details, routes, timings, and fares. The application is developed using React.js to provide a smooth and user-friendly experience.

\section{Study of Existing Bus Booking Platforms}

Several popular online bus booking platforms such as RedBus, AbhiBus, and MakeMyTrip are widely used by passengers across the country. These platforms provide features like real-time seat availability, route selection, bus timing details, and fare comparison. They also support multiple payment options including debit cards, credit cards, UPI, and net banking. Due to these features, they have become reliable solutions for online bus ticket booking.

Although these platforms are feature-rich, they often suffer from overly complex user interfaces. New users may find it difficult to navigate through multiple filters, pop-ups, and advertisements displayed on the screen. The presence of excessive promotional content can distract users from the main booking process. This complexity may reduce user satisfaction, especially for first-time users.

Another drawback observed in existing systems is the lack of customization and clarity in seat selection and passenger details entry. Some platforms require users to go through many steps before completing a booking. This can increase the time taken to book tickets and may lead to confusion. Additionally, unnecessary options and recommendations can overwhelm users during the booking flow.

These limitations influenced the design of the proposed bus booking system. The focus was placed on creating a clean, simple, and user-friendly interface with minimal distractions. Only essential features required for booking were included to improve usability. The aim is to provide a smooth booking experience while maintaining all necessary functionalities.

\section*{Comparative Analysis of Existing Systems}
\addcontentsline{toc}{section}{Comparative Analysis of Existing Systems}

A comparative analysis of existing online bus booking systems highlights both their advantages and shortcomings. Most current platforms offer essential features such as bus search, seat selection, route details, and online payment facilities. However, many of these systems suffer from complex user interfaces, excessive advertisements, and limited administrative control, which negatively impact the overall user experience. In several cases, customization options for managing buses, fares, and schedules are restricted, making system maintenance less flexible.

The proposed system is designed to overcome these limitations by providing a clean, simple, and user-friendly interface for passengers, along with a dedicated admin dashboard for effective management. Emphasis is placed on responsive design, clear navigation, and modular components to ensure smooth usability across devices. The system allows administrators to manage bus details, routes, schedules, and pricing efficiently. This comparative study helped identify key gaps in existing solutions, which were addressed during the design and implementation of the proposed bus booking system.
